\documentclass[12pt, a4paper]{report}
\usepackage[utf-8]{inputenc}
\usepackage[vietnamese]{babel}
\usepackage{xcolor}
\usepackage{listings}
\usepackage{geometry}
\usepackage{fancyhdr}
\usepackage{hyperref}
\usepackage{tikz}
\usepackage{tcolorbox}
\usepackage{array}
\usepackage{booktabs}

\geometry{margin=2.5cm}

\lstset{
    basicstyle=\ttfamily\small,
    breaklines=true,
    columns=fullflexible,
    backgroundcolor=\color{gray!10},
    frame=single,
    language=Python
}

\pagestyle{fancy}
\fancyhf{}
\rhead{\thepage}
\lhead{Task Generation System}

\title{\textbf{GIẢI THÍCH CHI TIẾT HỆ THỐNG TASK GENERATION}\\
\Large AI-Powered Requirement Analysis \& Task Generator}

\author{AI Project - Requirement Analyzer}
\date{\today}

\begin{document}

\maketitle
\tableofcontents
\newpage

\chapter{GIỚI THIỆU HỆ THỐNG}

\section{Task Generation System là gì?}

\textbf{Task Generation System} là một hệ thống AI powerpoint được xây dựng để tự động **chuyển đổi tài liệu yêu cầu (Requirements Document) thành các Task/User Story có cấu trúc**.

\begin{tcolorbox}[colback=blue!5, colframe=blue!50]
\textbf{Mục đích chính:}
\begin{itemize}
    \item 🎯 Tự động phân tích tài liệu yêu cầu
    \item 🎯 Tách yêu cầu từ tài liệu thô
    \item 🎯 Tạo Task/User Story theo chuẩn Agile
    \item 🎯 Phân loại, ước tính Story Points
    \item 🎯 Sẵn sàng upload lên Jira
\end{itemize}
\end{tcolorbox}

\section{Quy Trình Xử Lý - Pipeline}

Hệ thống hoạt động theo một pipeline 5 bước:

\begin{tcolorbox}[colback=yellow!5, colframe=yellow!50]
\begin{verbatim}
[RAW DOCUMENT]
       │
       ↓
[STAGE 1: SEGMENTATION]
Tách tài liệu thành các câu/đoạn
       │
       ↓
[STAGE 2: REQUIREMENT DETECTION]
Phát hiện câu nào là yêu cầu thực tế
       │
       ↓
[STAGE 3: ENRICHMENT]
Phân loại: Type, Priority, Domain, Role
       │
       ↓
[STAGE 4: TASK GENERATION]
Tạo Task/User Story chi tiết
       │
       ↓
[STAGE 5: POST-PROCESSING]
Loại bỏ trùng lặp, định dạng lại
       │
       ↓
[STRUCTURED TASKS OUTPUT]
(JSON - Ready for Jira)
\end{verbatim}
\end{tcolorbox}

\newpage

\chapter{CÁC STAGE CHI TIẾT}

\section{Stage 1: Segmentation (Tách Tài Liệu)}

\subsection{Mục Đích}

Chia tài liệu yêu cầu thành các đơn vị xử lý nhỏ (câu, đoạn).

\subsection{Đầu Vào \& Đầu Ra}

\begin{table}[h]
\centering
\begin{tabular}{|p{3cm}|p{6cm}|}
\hline
\textbf{Input} & \textbf{Output} \\
\hline
Raw document (TXT, MD, DOCX, PDF) & 
List of Sentences with metadata:
\begin{itemize}
\item Text content
\item Line number
\item Doc offset
\item Section header
\end{itemize} \\
\hline
\end{tabular}
\end{table}

\subsection{Ví Dụ}

\begin{tcolorbox}[colback=cyan!5]
\textbf{Input Document:}
\begin{lstlisting}
Requirements for Hotel Management System:

1. The system must allow customers 
   to search for available rooms.
2. Check-in and check-out dates 
   are required.
\end{lstlisting}

\textbf{Output (Segmented):}
\begin{lstlisting}
Sentences:
[0] "The system must allow customers 
     to search for available rooms."
    - line: 3
    - section: "Requirements for Hotel 
                Management System"
    
[1] "Check-in and check-out dates 
    are required."
    - line: 5
    - section: "Requirements for Hotel 
                Management System"
\end{lstlisting}
\end{tcolorbox}

\section{Stage 2: Requirement Detection (Phát Hiện)}

\subsection{Mục Đích}

Phát hiện **which sentences are actual requirements** và **filter out noise** (tiêu đề, giải thích, v.v.).

\subsection{Cách Hoạt Động}

Sử dụng **trained ML classifier** để dự đoán xác suất một câu là yêu cầu:

\begin{tcolorbox}[colback=orange!5]
\textbf{Requirement Detector Model:}
\begin{itemize}
    \item \textbf{Input:} Câu text
    \item \textbf{Feature Extraction:} TF-IDF vectorization
    \item \textbf{Model:} Logistic Regression / SVM
    \item \textbf{Output:} Probability [0-1]
    \item \textbf{Threshold:} > 0.5 = Yêu cầu (tunable)
\end{itemize}

\textbf{Ví dụ dự đoán:}
\begin{lstlisting}
"The system must allow customers to search" 
→ Probability: 0.92 ✓ (IS Requirement)

"Hotel Management System" 
→ Probability: 0.15 ✗ (NOT Requirement)

"Customers can modify their bookings" 
→ Probability: 0.87 ✓ (IS Requirement)
\end{lstlisting}
\end{tcolorbox}

\section{Stage 3: Enrichment (Phân Loại)}

\subsection{Mục Đích}

Gắn **metadata vào mỗi requirement**: Type, Priority, Domain, Role.

\subsection{Multi-Class Classifiers}

\begin{table}[h]
\centering
\begin{tabular}{|p{2.5cm}|p{4cm}|p{3cm}|}
\hline
\textbf{Label} & \textbf{Classes} & \textbf{Ví Dụ} \\
\hline
\textbf{Type} & Functional, Security, Interface, Data, Performance & "User can login" → Functional \\
\hline
\textbf{Priority} & Critical, High, Medium, Low & "Must encrypt data" → Critical \\
\hline
\textbf{Domain} & ecommerce, iot, healthcare, finance, general & "Payment processing" → Finance \\
\hline
\textbf{Role} & Backend, Frontend, QA, DevOps, BA, Security & "Database setup" → DevOps \\
\hline
\end{tabular}
\end{table}

\subsection{Ví Dụ Enrichment}

\begin{tcolorbox}[colback=green!5]
\textbf{Input (Requirement):}
\begin{lstlisting}
"The system must encrypt customer 
payment information in transit and at rest"
\end{lstlisting}

\textbf{Output (Enriched):}
\begin{lstlisting}
GeneratedTask {
    title: "Encrypt customer payment data",
    type: "security",           ← Classifier 1
    priority: "Critical",       ← Classifier 2
    domain: "finance",          ← Classifier 3
    role: "Backend",            ← Classifier 4
    ...
}
\end{lstlisting}
\end{tcolorbox}

\section{Stage 4: Task Generation (Tạo Task)}

\subsection{3 Modes Tạo Task}

Hệ thống hỗ trợ **3 cách khác nhau** để tạo task:

\subsubsection{Mode 1: Template-Based (Nhanh \& Consistent)}

\begin{tcolorbox}[colback=red!5, colframe=red!50]
\textbf{Đặc điểm:}
\begin{itemize}
    \item ⚡ Nhanh nhất (miliseconds)
    \item 📋 Sử dụng templates cố định
    \item 🎯 Output consistent
    \item 💾 Không cần API
\end{itemize}

\textbf{Cách hoạt động:}
\begin{lstlisting}
1. Extract entities từ requirement 
   (subject, action, object, condition)
   
2. Match với template tương ứng
   
3. Fill-in template với entities

Template ví dụ (Functional):
Title: "Implement {action} for {object}"
Description: "The system needs to {action} 
              {object} {condition}."
AC: - User can {action} {object}
    - System validates input
    - Error handling
\end{lstlisting}
\end{tcolorbox}

\subsubsection{Mode 2: LLM-Based (Linh Hoạt \& Natural)}

\begin{tcolorbox}[colback=purple!5, colframe=purple!50]
\textbf{Đặc điểm:}
\begin{itemize}
    \item 🤖 Sử dụng OpenAI/Anthropic API
    \item 📝 Natural language output
    \item 🎯 Context-aware (hiểu ngữ cảnh)
    \item 💰 Tốn chi phí API
    \item ⏱️ Chậm hơn template
\end{itemize}

\textbf{Cách hoạt động:}
\begin{lstlisting}
LLMTaskGenerator
   ↓
1. Gửi requirement đến LLM (GPT-4, Claude)
   với prompt structured
   
2. LLM tạo task output theo schema JSON
   
3. Validate response
   
4. Return GeneratedTask

Prompt ví dụ:
"Given this requirement: {requirement}
Generate a Jira-like task with:
- Title (action-oriented)
- Description
- 3-5 Acceptance Criteria
Output as JSON matching this schema: {...}"
\end{lstlisting}
\end{tcolorbox}

\subsubsection{Mode 3: Model-Based (Production - No APIs)}

\begin{tcolorbox}[colback=teal!5, colframe=teal!50]
\textbf{Đặc điểm:}
\begin{itemize}
    \item 🚀 Production ready
    \item 💾 Không cần API keys
    \item ⚡ Nhanh \& offline
    \item 🎯 Trained on project data
    \item 🔒 Data privacy (on-device)
\end{itemize}

\textbf{Cách hoạt động:}
\begin{lstlisting}
ModelBasedTaskGenerator
   ↓
1. Load trained ML models
   (requirement detector, enrichers)
   
2. Extract entities using spaCy NLP
   
3. Apply generation patterns
   (NOT hardcoded templates)
   
4. Generate natural variations
   using trained models
   
5. Return GeneratedTask

Models trained on:
- Project requirements dataset
- Historical tasks
- Requirement patterns
\end{lstlisting}
\end{tcolorbox}

\subsection{So Sánh 3 Modes}

\begin{table}[h]
\centering
\begin{tabular}{|l|c|c|c|c|}
\hline
\textbf{Aspect} & \textbf{Template} & \textbf{LLM} & \textbf{Model} \\
\hline
Speed & ⚡⚡⚡ & ⚡ & ⚡⚡⚡ \\
Flexibility & ⭐⭐ & ⭐⭐⭐ & ⭐⭐⭐ \\
Cost & Free & \$\$\$ & Free \\
API needed & No & Yes & No \\
Training data & Hardcoded & In-context & External \\
Natural output & ⭐⭐ & ⭐⭐⭐ & ⭐⭐⭐ \\
\hline
\end{tabular}
\end{table}

\section{Stage 5: Post-Processing}

\subsection{Các Bước}

\begin{enumerate}
    \item \textbf{Deduplication:} Loại bỏ task trùng lặp
    \item \textbf{Validation:} Kiểm tra required fields
    \item \textbf{Normalization:} Chuẩn hóa format
    \item \textbf{Filtering:} Loại bỏ low-confidence tasks
    \item \textbf{Sorting:} Sắp xếp theo priority
\end{enumerate}

\newpage

\chapter{TASK SCHEMA - CỬ TRÚC DỮ LIỆU}

\section{GeneratedTask Model}

Mỗi task sinh ra có cấu trúc Pydantic schema như sau:

\begin{tcolorbox}[colback=gray!5]
\begin{lstlisting}
class GeneratedTask(BaseModel):
    # ID & Identifiers
    task_id: str = UUID
    epic: Optional[str]
    module: Optional[str]
    
    # Core Content
    title: str                    # "Implement user login"
    description: str              # Detailed paragraph
    acceptance_criteria: List[str] # 3-5 AC items
    
    # Classification Labels
    type: str          # functional, security, interface...
    priority: str      # Critical, High, Medium, Low
    domain: str        # ecommerce, healthcare, finance...
    role: str          # Backend, Frontend, QA, DevOps...
    labels: List[str]  # Additional tags
    
    # Estimation
    story_points: Optional[int]    # 1, 2, 3, 5, 8, 13, 21
    estimated_hours: Optional[float]
    confidence: float              # [0.0 - 1.0]
    complexity: Optional[str]      # Simple, Medium, Complex
    
    # Source Tracking
    source: Optional[TaskSource]
    - sentence: str        # Original requirement
    - section: Optional[str]
    - doc_offset: [start, end]
    - line_number: int
    
    # Dependencies
    dependencies: List[str] # task IDs this depends on
    
    # Metadata
    generated_at: datetime
    generator_version: str
\end{lstlisting}
\end{tcolorbox}

\newpage

\chapter{PIPELINE CODE - HOW IT WORKS}

\section{Main Pipeline Class}

\begin{tcolorbox}[colback=blue!5]
\begin{lstlisting}
class TaskGenerationPipeline:
    """
    Main orchestrator combining all stages
    """
    
    def __init__(
        self,
        generator_mode: str = "template",
        # "template" (fast)
        # "llm" (flexible, needs API)
        # "model" (production, trained)
        llm_provider: Optional[str] = None,
        llm_model: Optional[str] = None,
        llm_api_key: Optional[str] = None
    ):
        self.mode = generator_mode
        
        # Initialize components
        self.segmenter = get_segmenter()
        self.detector = get_detector()    # Req detection
        self.enricher = get_enrichment_pipeline()
        
        # Choose generator based on mode
        if mode == "template":
            self.generator = TemplateTaskGenerator()
        elif mode == "llm":
            self.generator = LLMTaskGenerator(
                provider=llm_provider,
                model=llm_model,
                api_key=llm_api_key
            )
        elif mode == "model":
            self.generator = ModelBasedTaskGenerator()
    
    def generate_tasks(
        self,
        document: str,
        max_tasks: int = 100,
        min_confidence: float = 0.5
    ) -> TaskGenerationResponse:
        """Main execution method"""
        
        # Stage 1: Segment
        sentences = self.segmenter.segment(document)
        logger.info(f"✓ Segmented into {len(sentences)} sentences")
        
        # Stage 2: Detect requirements
        requirements = []
        for sent in sentences:
            prob = self.detector.predict_proba(sent.text)
            if prob > min_confidence:
                requirements.append(sent)
        logger.info(f"✓ Detected {len(requirements)} requirements")
        
        # Stage 3: Enrich
        enriched = []
        for req in requirements:
            enriched_req = self.enricher.enrich(req)
            enriched.append(enriched_req)
        logger.info(f"✓ Enriched with metadata")
        
        # Stage 4: Generate tasks
        tasks = []
        for req in enriched[:max_tasks]:
            task = self.generator.generate(req)
            tasks.append(task)
        logger.info(f"✓ Generated {len(tasks)} tasks")
        
        # Stage 5: Post-process
        tasks = self.postprocessor.process(tasks)
        logger.info(f"✓ Post-processed to {len(tasks)} tasks")
        
        return TaskGenerationResponse(
            tasks=tasks,
            total_count=len(tasks),
            processing_time_ms=...
        )
\end{lstlisting}
\end{tcolorbox}

\section{Usage Example}

\begin{tcolorbox}[colback=green!5]
\begin{lstlisting}
from requirement_analyzer.task_gen import get_pipeline

# Initialize with different modes
pipeline_fast = get_pipeline(mode="template")
pipeline_flex = get_pipeline(
    mode="llm",
    llm_provider="openai",
    llm_model="gpt-4o-mini"
)
pipeline_prod = get_pipeline(mode="model")

# Use pipeline
document = """
Hotel Management System Requirements:

1. The system must allow customers to search
   for available rooms by date range.
2. Customers can view room details including
   price, amenities, and availability.
3. System must send confirmation email
   after booking is confirmed.
4. Payment must be encrypted using SSL/TLS.
"""

# Generate tasks
result = pipeline_fast.generate_tasks(
    document,
    max_tasks=20,
    min_confidence=0.6
)

# Access results
print(f"Generated {result.total_count} tasks")
for task in result.tasks:
    print(f"\n{task.title}")
    print(f"  Priority: {task.priority}")
    print(f"  Story Points: {task.story_points}")
    print(f"  Type: {task.type}")
    print(f"  Role: {task.role}")
    print(f"  Confidence: {task.confidence:.2%}")
    print(f"  AC: {len(task.acceptance_criteria)}")

# Export to Jira or JSON
result.export_to_jira(project_key="HMS")
result.export_to_json("generated_tasks.json")
\end{lstlisting}
\end{tcolorbox}

\newpage

\chapter{USER STORIES VÀ TASKS TRONG HỆ THỐNG}

\section{Task Generation Feature - User Stories}

\subsection{US1: Tạo Requirement Document}

\begin{tcolorbox}[colback=cyan!5, colframe=cyan!50]
\begin{lstlisting}
Title: Upload requirement document

As a Project Manager
I want to upload requirement documents
So that I can generate tasks automatically

Acceptance Criteria:
✓ User can upload TXT, MD, DOCX, PDF files
✓ System validates file format
✓ Shows upload progress
✓ Displays file content preview
✓ Error handling for invalid files

Type: Functional
Priority: High
Story Points: 3
\end{lstlisting}
\end{tcolorbox}

\subsection{US2: Generate Tasks from Requirements}

\begin{tcolorbox}[colback=orange!5, colframe=orange!50]
\begin{lstlisting}
Title: Generate tasks with multiple modes

As a Project Manager
I want to generate tasks using multiple 
generation modes (Template/LLM/Model)
So that I can choose best approach for 
my project needs

Acceptance Criteria:

Scenario 1: Template Mode
✓ Select "Fast Template Generation"
✓ System uses predefined templates
✓ Tasks generated within 2 seconds
✓ Output is consistent across runs

Scenario 2: LLM Mode
✓ Select "Flexible LLM Generation"
✓ Configure LLM provider (OpenAI/Anthropic)
✓ System generates natural task descriptions
✓ Request/response logged for audit

Scenario 3: Model Mode
✓ Select "Production Model"
✓ System uses locally trained models
✓ No API keys required
✓ Fast processing without internet

Type: Functional
Priority: High
Story Points: 8
Complexity: Complex
\end{lstlisting}
\end{tcolorbox}

\subsection{US3: Review \& Edit Generated Tasks}

\begin{tcolorbox}[colback=yellow!5, colframe=yellow!50]
\begin{lstlisting}
Title: Review and edit generated tasks

As a Project Manager
I want to review, edit, or delete 
generated tasks
So that I can ensure quality before 
importing to Jira

Acceptance Criteria:
✓ Display tasks in table/list view
✓ Can edit title, description, AC
✓ Can change priority/type/story points
✓ Can delete unwanted tasks
✓ Show confidence score for each task
✓ Bulk operations (select multiple)
✓ Changes saved to local storage

Type: Interface
Priority: High
Story Points: 5
\end{lstlisting}
\end{tcolorbox}

\subsection{US4: Export Tasks to Jira}

\begin{tcolorbox}[colback=purple!5, colframe=purple!50]
\begin{lstlisting}
Title: Export generated tasks to Jira

As a Project Manager
I want to export tasks to Jira directly
So that I can quickly populate project board

Acceptance Criteria:
✓ Authenticate with Jira API
✓ Select target Jira project
✓ Map task fields to Jira fields
✓ Support custom fields
✓ Show import progress
✓ Handle duplicate prevention
✓ Create tasks in batches
✓ Error logging for failed imports

Type: Integration
Priority: High
Story Points: 8
Dependencies: [US3, Jira API setup]
\end{lstlisting}
\end{tcolorbox}

\section{Related Backend Tasks}

\subsection{Task: Implement Requirement Detector ML Model}

\begin{tcolorbox}[colback=red!5]
\begin{lstlisting}
Title: Train requirement detection classifier

Description:
Develop and train ML model to classify
whether a sentence is an actual requirement

Technical Details:
- Model type: Logistic Regression / SVM
- Features: TF-IDF vectors
- Training data: 1000+ labeled sentences
- Classes: [Requirement, Non-requirement]
- Target accuracy: >90%

Acceptance Criteria:
✓ Model trained on dataset
✓ Cross-validation accuracy >90%
✓ Model exported and serialized
✓ Inference time <10ms per sentence
✓ Model integrated with pipeline
✓ Unit tests passing

Type: Data/ML
Priority: Critical
Story Points: 13
Complexity: Complex
\end{lstlisting}
\end{tcolorbox}

\subsection{Task: Implement Enrichment Classifiers}

\begin{tcolorbox}[colback=magenta!5]
\begin{lstlisting}
Title: Train multi-class enrichment classifiers

Description:
Build 4 separate classifiers for:
1. Type classifier (7 classes)
   functional, security, interface, data, 
   performance, integration, documentation
   
2. Priority classifier (4 classes)
   Critical, High, Medium, Low
   
3. Domain classifier (6 classes)
   ecommerce, iot, healthcare, finance, 
   education, general
   
4. Role classifier (6 classes)
   Backend, Frontend, QA, DevOps, BA, Security

Acceptance Criteria:
✓ Each classifier accuracy >85%
✓ Fast inference (<100ms per sentence)
✓ Models serialized and versioned
✓ Integrated in enrichment pipeline
✓ Fallback to defaults if confidence low
✓ Pipeline tests passing

Type: Data/ML
Priority: High
Story Points: 21
Complexity: Very Complex
\end{lstlisting}
\end{tcolorbox}

\subsection{Task: Implement LLM Integration}

\begin{tcolorbox}[colback=blue!5]
\begin{lstlisting}
Title: Integrate with LLM API (OpenAI/Anthropic)

Description:
Implement LLMTaskGenerator class to use
external LLM APIs for task generation

Technical Requirements:
- Support OpenAI (GPT-4, GPT-4o-mini)
- Support Anthropic (Claude-3-Haiku)
- Structured JSON output
- Streaming support
- Error handling \& retries
- Cost tracking (API usage)
- Rate limiting

Acceptance Criteria:
✓ OpenAI integration working
✓ Anthropic integration working
✓ Structured prompt engineering
✓ JSON validation
✓ Caching implemented
✓ Error messages clear
✓ Cost calculation accurate
✓ Integration tests passing

Type: Integration
Priority: High
Story Points: 8
\end{lstlisting}
\end{tcolorbox}

\subsection{Task: Implement Model-Based Generator}

\begin{tcolorbox}[colback=green!5]
\begin{lstlisting}
Title: Implement production model-based generator

Description:
Build ModelBasedTaskGenerator using trained
ML models for on-device task generation

Technical Details:
- Load trained models from disk
- Extract entities with spaCy NLP
- Generate task variations
- No external API required

Acceptance Criteria:
✓ Models loaded correctly
✓ Entity extraction working
✓ Task generation deterministic
✓ Performance >1000 tasks/min
✓ Memory usage <500MB
✓ All AC met
✓ E2E tests passing
✓ Documentation complete

Type: Backend
Priority: High
Story Points: 13
\end{lstlisting}
\end{tcolorbox}

\newpage

\chapter{KIẾN TRÚC \& TECHNICAL DETAILS}

\section{System Architecture}

\begin{tcolorbox}[colback=gray!5]
\begin{verbatim}
┌─────────────────────────────────────────┐
│  FastAPI Web Server                     │
│  (routes/api endpoints)                 │
└──────────────────┬──────────────────────┘
                   │
    ┌──────────────┴──────────────┐
    │                             │
┌───▼────────────────┐  ┌────────▼──────┐
│ Task Gen Pipeline  │  │ File Upload   │
│                    │  │ Processing    │
└───┬────────────────┘  └────────▲──────┘
    │                            │
    ├─── Segmenter             
    │    (sentence splitting)
    │
    ├─── Detector              
    │    (ML classifier)
    │    models/requirement_detector_model.joblib
    │
    ├─── Enricher              
    │    (4 classifiers)
    │    - type_classifier
    │    - priority_classifier
    │    - domain_classifier
    │    - role_classifier
    │
    ├─── Generator (3 modes)       
    │    ├─ TemplateGenerator (fast)
    │    ├─ LLMGenerator (flexible)
    │    └─ ModelGenerator (production)
    │
    └─── PostProcessor          
         (dedup, validation)
         │
         └─── Output: JSON/Jira

┌─────────────────────────────────────────┐
│  Database (PostgreSQL)                  │
│  - tasks table                          │
│  - document_uploads                     │
│  - enrichment_logs                      │
└─────────────────────────────────────────┘
\end{verbatim}
\end{tcolorbox}

\section{Component Dependencies}

\begin{table}[h]
\centering
\begin{tabular}{|l|p{5cm}|p{3cm}|}
\hline
\textbf{Component} & \textbf{Dependencies} & \textbf{Purpose} \\
\hline
Segmenter & nltk, spaCy & Sentence tokenization \\
\hline
Detector & scikit-learn, joblib & ML classification \\
\hline
Enrichers & scikit-learn, joblib & Multi-class labeling \\
\hline
Template Gen & spaCy, regex & Pattern matching \\
\hline
LLM Gen & openai/anthropic & External API calls \\
\hline
Model Gen & spaCy, joblib & On-device inference \\
\hline
\end{tabular}
\end{table}

\section{Data Flow Example}

\begin{tcolorbox}[colback=yellow!5]
\begin{lstlisting}
EXAMPLE: User uploads hotel_requirements.md

1. UPLOAD STAGE
   File: hotel_requirements.md
   Size: 50 KB
   Status: Uploaded ✓

2. SEGMENTATION
   Input: Raw text (1200 words)
   Process: Split into sentences
   Output: 98 sentences with metadata

3. REQUIREMENT DETECTION
   Processed: 98 sentences
   Detected: 74 are requirements (>0.5 confidence)
   Filtered out: 24 non-requirement sentences

4. ENRICHMENT (74 requirements)
   Type Distribution:
   - Functional: 45 (60%)
   - Security: 12 (16%)
   - Data: 10 (13%)
   - Performance: 5 (7%)
   - Interface: 2 (3%)

   Priority Distribution:
   - Critical: 8
   - High: 34
   - Medium: 28
   - Low: 4

5. TASK GENERATION (Mode: Template)
   Time: 280ms
   Tasks generated: 74
   Avg confidence: 0.85

6. POST-PROCESSING
   Removed duplicates: 3
   Final tasks: 71
   Ready for Jira export

7. OUTPUT
   Format: JSON
   Size: 125 KB
   Status: Ready for import
\end{lstlisting}
\end{tcolorbox}

\newpage

\chapter{CÓ GÌ VUI VẪN VỀ HỆ THỐNG}

\section{Advantages}

\begin{itemize}
    \item ✅ \textbf{Tự động hóa:} Không cần tạo task thủ công
    \item ✅ \textbf{Consistent:} Tất cả task cùng format, structure
    \item ✅ \textbf{Flexible:} 3 generation modes với trade-offs khác nhau
    \item ✅ \textbf{Smart:} ML classifiers learn từ dữ liệu
    \item ✅ \textbf{Jira-ready:} Export directly to Jira/JSON
    \item ✅ \textbf{Scalable:} Handle 1000+ requirements
    \item ✅ \textbf{Traceable:} Track source sentence cho mỗi task
\end{itemize}

\section{Limitations \& Future Work}

\begin{itemize}
    \item ⚠️ Không perfect 100\% - cần human review
    \item ⚠️ Language-specific - models trained for English
    \item ⚠️ Requires clean input - garbage in = garbage out
    \item ⚠️ ML models need retraining với new patterns
    
    \item 🔮 Future: Multi-language support
    \item 🔮 Future: Custom model training per-project
    \item 🔮 Future: Story point auto-estimation (ML)
    \item 🔮 Future: Dependency detection between tasks
    \item 🔮 Future: Real-time feedback \& refinement
\end{itemize}

\chapter{SUMMARY}

\begin{tcolorbox}[colback=purple!5, colframe=purple!50]
\textbf{Task Generation System} của bạn là một **production-grade AI system** sau:

\begin{enumerate}
    \item 📄 **Input:** Raw requirement documents (TXT, MD, DOCX, PDF)
    \item 🔄 **Process:** 5-stage pipeline (Segment → Detect → Enrich → Generate → Post-process)
    \item 🤖 **3 Generation Modes:}
    \begin{itemize}
        \item Template (fast, consistent, no API)
        \item LLM (flexible, natural, needs API)
        \item Model (production, trained, no API)
    \end{itemize}
    \item 📊 **Output:** Structured tasks (JSON/Jira)
    \item 🎯 **Benefits:** Automation, consistency, scalability, quality
\end{enumerate}

\textbf{Cách sử dụng:}
\begin{lstlisting}
# Fast generation
tasks = pipeline.generate_tasks(doc, mode="template")

# Flexible generation  
tasks = pipeline.generate_tasks(doc, mode="llm", 
                                 provider="openai")

# Production
tasks = pipeline.generate_tasks(doc, mode="model")
\end{lstlisting}
\end{tcolorbox}

\end{document}
